\section{问卷质量分析}\label{sec:quality}

\subsection{随机性检验}

为检验样本是否存在抽样偏差，本研究采用\textbf{游程检验}，利用SPSS对性别、年龄、现居住地、现居住区域、职业类型、文化程度6项人口统计学变量进行分析。该检验通过比较实际游程数与随机条件下的期望游程数，判断样本在采集顺序上的随机程度；检验值按各变量的中位数设定（性别取切分值2），有效样本为467份，结果如表\ref{tab:runs}所示。
\begin{table}[H]
\centering
\caption{游程检验结果}
\label{tab:runs}
\rowcolors{2}{tablight}{white}
\begin{tabular}{M{2.8cm}M{1.6cm}M{1.6cm}M{1.6cm}M{1.6cm}M{1.6cm}M{1.6cm}}
\bluetoprule
\rowcolor{tabhead}
\thead{参数} & \thead{性别} & \thead{年龄} & \thead{现居住地} & \thead{现居住区域} & \thead{职业类型} & \thead{文化程度} \\
\hline
检验值 & 2 & 3 & 2 & 2 & 4 & 4 \\
个案数$<$检验值 & 260 & 167 & 132 & 134 & 196 & 195 \\
个案数$\ge$检验值 & 207 & 300 & 335 & 333 & 271 & 272 \\
总个案数 & 467 & 467 & 467 & 467 & 467 & 467 \\
游程数 & 177 & 174 & 157 & 164 & 204 & 197 \\
$Z$ & -5.115 & -4.191 & -3.815 & -3.183 & -2.328 & -2.967 \\
渐近显著性(双尾) & $<0.001$ & $<0.001$ & $<0.001$ & 0.001 & 0.020 & 0.003 \\
\bluebottomrule
\end{tabular}
\end{table}

检验结果显示，6项变量的显著性均小于0.05，说明样本在采集顺序上并非完全随机。其中，性别、年龄和现居住地的显著性均小于0.001，职业类型的显著性最高但也仅为0.020，表明各变量均存在不同程度的系统性排列。该现象主要与问卷按城市分批发放和回收有关，同一城市样本在数据序列中相对集中；其反映的是数据采集顺序特征，并不影响样本的覆盖范围与代表性，数据仍可用于后续分析。

\subsection{相关性分析}

由于“受噪声影响程度”采用反向计分，分析前先进行正向化处理；考虑到相关变量多为有序变量，本研究采用\textbf{Spearman秩相关分析}，利用SPSS对467份有效样本进行检验，以识别居民整体噪声感知与各类噪声源、噪声负面影响之间的关联，结果如图\ref{fig:spearman}所示。

\fig[0.98]{spearman_correlation.png}{居民整体噪声感知与各题项的Spearman相关分析结果\label{fig:spearman}}

结果显示，各题项与整体噪声感知\textbf{均呈显著正相关}（$P<0.05$），多数达到$P<0.001$水平。其中，交通噪声影响、公共活动噪声影响和装修/施工噪声影响的关联较强，是影响整体声环境体验的主要噪声源；其余题项多为弱相关，但显著性均满足要求。整体而言，各类噪声干扰越强，居民的睡眠、注意力、情绪和身心健康所受负面影响越明显。

\subsection{信度分析}

为检验问卷量表的可靠性与内部一致性，本研究采用\textbf{克朗巴哈系数（Cronbach's Alpha）}对22个量表题项进行信度分析（含14个噪声影响题项与受噪声影响程度、隔音水平评价、离噪声源距离、两项满意度评价及三项治理认知等8个其他量表题项）。其中“受噪声影响程度”为\textbf{反向计分}题（1=影响非常大，5=几乎没有影响），计算前已作正向化处理，结果如下表。

\begin{table}[H]
\centering
\caption{信度分析结果}
\rowcolors{2}{tablight}{white}
\begin{tabular}{M{8cm}M{5cm}}
\bluetoprule
\rowcolor{tabhead}
\thead{指标} & \thead{数值} \\
\hline
克隆巴赫Alpha & 0.735 \\
项数 & 22 \\
\bluebottomrule
\end{tabular}
\end{table}

结果显示，克隆巴赫Alpha为0.735，表明问卷量表具有可接受的信度。

\subsection{效度检验}

为检验量表题项的结构效度，本研究采用\textbf{KMO取样适切性量数}和\textbf{巴特利特球形检验}，利用SPSS对参与后续因子分析的14个噪声影响题项进行分析，结果如下表。

\begin{table}[H]
\centering
\caption{效度检验结果}
\rowcolors{2}{tablight}{white}
\begin{tabular}{M{8cm}M{5cm}}
\bluetoprule
\rowcolor{tabhead}
\thead{指标} & \thead{数值} \\
\hline
KMO取样适切性量数 & 0.815 \\
巴特利特球形检验近似卡方 & 2119.289 \\
自由度 & 91 \\
显著性 & $<0.001$ \\
\bluebottomrule
\end{tabular}
\end{table}

结果显示，KMO值为0.815，巴特利特球形检验显著性小于0.001，表明量表适合进行因子分析，结构效度较好。
